\section{Related Work}

\noindent\textbf{Spatio-Temporal Imputation under Sparse Observability.}
Classical methods such as Kriging \citep{cressie1990origins} and statistical models with spatio-temporal kernels \citep{wikle2019spatio,de2005predictive} provide principled approaches for interpolation from sparse measurements, but they typically rely on stationarity, smoothness, and kernel assumptions that may not hold across heterogeneous physical systems. They also scale poorly, with cubic complexity in the number of observations, and do not explicitly account for governing dynamics during interpolation. Physics-based numerical solvers such as finite difference and finite element methods \citep{leveque2007finite,hughes2003finite} provide mechanistic structure, but require sufficiently specified physical models, boundary conditions, and parameters, which are often unavailable or non-identifiable from sparse sensor data alone.

Learning-based approaches have been developed to model more complex spatio-temporal dependencies. Message-passing recurrent neural networks (MPRNNs) \citep{iyer2022modeling} use learned dynamics over sensor networks; graph-based approaches such as Graph WaveNet \citep{wu2019graph} and ST-Transformer \citep{xu2020spatial} encode spatial relations through fixed or learned adjacency structures; and transformer models such as ImputeFormer \citep{nie2024imputeformer} capture long-range temporal dependencies in structured multivariate time series. These methods are effective when the sensing topology and data-generating process align with their architectural priors, but they generally treat sparse sensing as a missing-data problem rather than as an observability-constrained physical inference problem. In contrast, our work explicitly studies reconstruction under severe sparse sensing, where the goal is not unconstrained global field recovery, but reliable sensor-space modeling and surrogate field construction under limited observational support.

\noindent\textbf{Mesh-Free vs.\ Mesh-Dependent Spatio-Temporal Modeling.}
Prior work on spatio-temporal reconstruction can be broadly divided into mesh-free and mesh-dependent approaches. Mesh-free methods~\citep{tancik2020fourier,sitzmann2020implicit,raissi2019physics,cressie1990origins,hensman2013gaussian} operate directly on continuous coordinates, support arbitrary query locations, and naturally accommodate irregular sensor layouts and multi-resolution inference. These models are attractive for sparse sensing because they are not tied to a fixed grid or sensor graph. However, under extreme sparsity, continuous query capability does not by itself guarantee identifiable global reconstruction: predictions away from observational support are necessarily governed by the model's implicit priors.

Mesh-dependent methods~\citep{nie2024imputeformer,iyer2022modeling,xu2020spatial,wu2019graph,liu2023physics,liang2024higher,santos2023development,zhao2024recfno,li2020fourier,rahman2022u,li2017diffusion} assume a fixed spatial discretization, such as dense grids or static sensor graphs, and typically treat missing values as masked entries rather than absent spatial observations. Such methods can learn strong global or low-rank priors over fixed domains, and may generalize well when these priors match the underlying data-generating process. However, they are less naturally suited to irregular, persistent sparse sensor networks where reliable prediction is concentrated around the deployed sensing topology. FieldFormer follows the mesh-free paradigm, but uses it for locality-aware sensor-space modeling rather than claiming unconstrained global recovery under sparse observability.

\noindent\textbf{Physics-Informed Neural Networks and Differentiable PDE Solvers.}
Physics-Informed Neural Networks (PINNs) \citep{raissi2019physics} embed PDE constraints into neural network training by computing differential operators on model outputs via automatic differentiation and penalizing residual violations. This enables both forward and inverse modeling when the governing equations and sufficient observations are available. However, PINNs typically rely on global coordinate function approximators and dense collocation over structured domains. As noted by \citep{krishnapriyan2021characterizing}, vanilla global PINNs often struggle in moderately complex PDE regimes and real-world scenarios. In sparse longitudinal sensing settings, where observations are dense in time but sparse in space, global PDE residual minimization may be weakly constrained or even misleading when parameters, forcing terms, and boundary conditions are unknown.

Other physics-integrated approaches, such as DiffTaichi \citep{hu2019difftaichi} and JAX-FDM \citep{kochkov2021machine}, enable differentiable programming over simulation pipelines, but still rely on mesh-based discretizations and sufficiently specified simulators. These assumptions limit their applicability when the available data consist of scattered sensor measurements and the underlying physical process is only partially known. FieldFormer instead incorporates physical structure as optional local regularization while prioritizing observational consistency over persistent sparse sensor networks.

\noindent\textbf{Transformers for Scientific Modeling and Field Inference.}
Transformer architectures have increasingly been applied to scientific modeling tasks, including spatio-temporal prediction and PDE-informed learning. TransFlowNet \citep{wang2022transflownet} introduces physics-constrained loss terms into transformer models for spatio-temporal flow prediction, while \citep{lorsung2024physics} integrate attention-based sequence modeling with physics-informed PDE tokens. These approaches typically assume dense spatial coverage, complete or well-specified governing physics, or grid-aligned discretizations, and often rely on global attention mechanisms. Such assumptions are difficult to satisfy in sparse sensing regimes with unobserved forcing, uncertain parameters, and incomplete boundary information.

In contrast, FieldFormer restricts attention to adaptive local neighborhoods while remaining globally grounded through continuous coordinates. This design reflects the view that, under sparse observability, reliable reconstruction depends on whether the sensor network sufficiently covers local domains of dependence. Rather than using transformers to model an unconstrained global latent field, FieldFormer uses local attention to organize sparse observations according to learned spatio-temporal geometry.

\noindent\textbf{Hybrid Approaches for Learning with Physical Structure.}
Recent work has explored hybrid approaches that combine learning with physical structure without explicitly solving governing PDEs. Neural Operators such as Fourier Neural Operators (FNOs) \citep{li2020fourier} and Graph Neural Operators \citep{li2020neural} learn mappings between function spaces, but typically require dense grid-based inputs and outputs. This makes them less suitable for scattered sensor measurements or persistent sparse sensing regimes where large portions of the domain are never observed. Other hybrids, including Physics-Informed Neural Fields \citep{chu2022physics} and Koopman operator methods \citep{lusch2018deep}, embed physical constraints or dynamical priors into latent representations, but generally assume dense trajectories, simulation-generated data, or sufficient observational coverage.

These approaches demonstrate the value of combining learning with physical structure, but they do not directly address the observability limits induced by sparse sensor networks. FieldFormer complements this line of work by focusing on locality-aware reconstruction under sparse observational support, where global field recovery is underconstrained and prediction quality depends critically on the alignment between model inductive bias and sensing topology.